\begin{table}[H]
  \centering
  \caption{Spending composition and residential sorting: Inc-bind cohort.}
  \label{tab:spending_decomposition}
  \small
  \begin{tabular}{lrrrr}
    \toprule
    Expenditure category & $N$ & Slope & SE & Corr. \\
    \midrule
    Debt service (exp.) & 19 & -0.629 & (0.577) & -0.256 \\
    EIT (rev. share) & 19 & -0.250 & (0.658) & -0.092 \\
    Police (exp.) & 39 & 0.194 & (0.713) & 0.045 \\
    Streets & roads (exp.) & 39 & 0.279 & (0.460) & 0.099 \\
    Intergovernmental (rev. share) & 19 & 0.478 & (0.606) & 0.188 \\
    Real estate (rev. share) & 19 & 0.543 & (0.925) & 0.141 \\
    \bottomrule
  \end{tabular}
  \vspace{0.3em}
  \begin{minipage}{0.95\textwidth}
    \footnotesize \textit{Notes:} For each Inc-bind
    municipality, the outcome is
    $\Delta\log(\text{live-and-work})$ measured from
    event-year $c{-}1$ to $c{+}5$. The covariate is
    $\Delta s^{X}$, the change in the share of
    total expenditure devoted to category $X$ over the
    same window. Slope is the OLS coefficient in the
    bivariate regression
    $\Delta\log N = \alpha + \gamma \, \Delta s^{X}$.
    A negative slope means the muni grew share-$X$ AND
    lost more residents; a positive slope means residents
    tolerated the hike when more of the new revenue went
    to category $X$. $N$ varies because some munis report
    zero (or no AFR detail) for some categories at one of
    the two event-window endpoints. $^{*}p<0.10$,
    $^{**}p<0.05$, $^{***}p<0.01$.
  \end{minipage}
\end{table}
